\documentclass[12pt]{article}
\usepackage[utf8]{inputenc}
\usepackage[T1]{fontenc}
\usepackage{hyperref}
\title{Mateq Language Resources for K-12 Curriculum Development}
\author{Compiled by Claude Sonnet 5.5 (claude-sonnet-5-5) for Andr\'as Kornai}
\date{}
\begin{document}
\maketitle

\section{Introduction}
Mateq (also Mate', Bumat\'eq, Kembayan; ISO 639-3 xem; Glottocode kemb1249) is an Austronesian (Dayak) language of Sanggau Regency, West Kalimantan, Indonesia. Estimates run from about 10,000 to 20,000 speakers; Wikipedia gives about 11,000 (2007). It should not be confused with Kemberano (bzp), a Papuan language of the Bird's Head. Searches found no school materials for Mateq.

\section{Primary Educational Resources}
No primers, textbooks or school programmes were found.

\section{Linguistic and Reference Materials}
\begin{itemize}
\item Timothy M. Connell (2013), MA thesis (University of Canterbury), a sketch grammar of Mateq. Glottolog record: \url{https://glottolog.org/resource/reference/id/442917}
\item Connell 2013 listing at CrossGram: \url{https://crossgram.clld.org/sources/bonmannsymmetrical-connell2013mateq}
\item Wikipedia, Kembayan language: \url{https://en.wikipedia.org/wiki/Kembayan_language}
\item Grambank entry: \url{https://grambank.clld.org/languages/kemb1249}
\item Omniglot, Kembayan (Mat\'eq) writing: \url{https://omniglot.com/writing/kembayan.htm}
\item Omniglot, Kembayan numbers: \url{https://www.omniglot.com/language/numbers/kembayan.htm}
\end{itemize}

\section{Religious and Community Resources}
Global Recordings Network lists a language page that appeared in the search results: \url{https://www.globalrecordings.net/en/language/5061}. The search did not show which materials it holds, so its contents are unverified.

\section{Digital and Audio Resources}
Apart from the Global Recordings page above, no audio, video or corpus resources were found.

\section{Curriculum Development Context}
Mateq has a sizeable speaker base and a published sketch grammar, but no known school materials. A curriculum could start from Connell's description and Omniglot's orthographic information, working with Indonesian local-content (muatan lokal) frameworks. A locally agreed orthography and a community materials team would be needed.
\end{document}